% ============================================================================
% MUSTERLÖSUNG: Verkehrssicherheit (LP-04) - Klasse 6
% Vollständiges AB mit eingesetzten Lösungen
% Stand: 04.02.2026
% ============================================================================

\documentclass[11pt,a4paper]{article}

\usepackage[utf8]{inputenc}
\usepackage[T1]{fontenc}
\usepackage[ngerman]{babel}
\usepackage{geometry}
\usepackage{helvet}
\usepackage{tikz}
\usetikzlibrary{decorations.pathreplacing}
\usepackage{fancyhdr}
\usepackage{xcolor}
\usepackage{amssymb}
\usepackage{tcolorbox}
\usepackage{qrcode}
\usepackage{array}
\usepackage{tabularx}
\usepackage[hidelinks]{hyperref}

\renewcommand{\familydefault}{\sfdefault}
\geometry{left=1.5cm, right=1.5cm, top=1.2cm, bottom=1.0cm}
\setlength{\parindent}{0pt}
\setlength{\parskip}{0pt}

% Farben
\definecolor{physik}{HTML}{2c5aa0}
\definecolor{reaktion}{HTML}{2196f3}
\definecolor{brems}{HTML}{ff9800}
\definecolor{anhalt}{HTML}{4caf50}
\definecolor{hinweis}{HTML}{fff3e0}
\definecolor{beispiel}{HTML}{e3f2fd}
\definecolor{loesung}{HTML}{c62828}

% Kopfzeile
\pagestyle{fancy}
\fancyhf{}
\fancyhead[L]{\small\textbf{Physik Klasse 6} | Verkehrssicherheit}
\fancyhead[R]{\small\textcolor{loesung}{\textbf{MUSTERLÖSUNG}}}
\renewcommand{\headrulewidth}{0.4pt}
\setlength{\headheight}{14pt}

% Boxen
\tcbuselibrary{skins}
\newtcolorbox{hinweisbox}[1][]{colback=hinweis, colframe=orange!70!black, boxrule=1pt, arc=3pt, left=6pt, right=6pt, top=4pt, bottom=4pt, #1}
\newtcolorbox{beispielbox}[1][]{colback=beispiel, colframe=physik, boxrule=1pt, arc=3pt, left=6pt, right=6pt, top=4pt, bottom=4pt, #1}
\newtcolorbox{merkbox}[1][]{colback=white, colframe=physik, boxrule=1pt, arc=3pt, left=6pt, right=6pt, top=4pt, bottom=4pt, #1}

% Befehle
\newcommand{\aufgabe}[1]{\textbf{\textcolor{physik}{#1}}}
\newcommand{\luecke}[1]{\underline{\hspace{#1}}}
\newcommand{\lsg}[1]{\textcolor{loesung}{\textbf{#1}}}

\begin{document}

% ============================================================================
% TITEL
% ============================================================================
\noindent
\begin{minipage}[t]{0.80\textwidth}
    \vspace{0pt}
    {\Large\bfseries Verkehrssicherheit: Reaktionsweg und Anhalteweg}\\[0.3em]
    Name: \lsg{Musterlösung} \hfill Datum: \lsg{04.02.2026}
\end{minipage}
\hfill
\begin{minipage}[t]{0.16\textwidth}
    \vspace{0pt}
    \raggedleft
    \href{https://devbox42.github.io/sciencesim/physik/kl06/02-bewegung/std-04-verkehrssicherheit/LP-04-verkehrssicherheit.html}{%
    \qrcode[height=1.4cm]{https://devbox42.github.io/sciencesim/physik/kl06/02-bewegung/std-04-verkehrssicherheit/LP-04-verkehrssicherheit.html}%
    }
\end{minipage}

\vspace{0.2em}
\begin{hinweisbox}
\centering
\textbf{Der Lernpfad führt dich Schritt für Schritt.} Scanne den QR-Code und fülle dieses Blatt parallel aus.
\end{hinweisbox}

\vspace{0.4em}

% ============================================================================
% LEGENDE
% ============================================================================
\begin{merkbox}
\textbf{Die drei wichtigen Wege:}
\hfill
\renewcommand{\arraystretch}{1.3}
\begin{tabular}{ll}
\textbf{Reaktionsweg} & = Weg während der Schrecksekunde (du fährst noch weiter!) \\
\textbf{Bremsweg} & = Weg während des Bremsens (bis das Auto steht) \\
\textbf{Anhalteweg} & = Reaktionsweg + Bremsweg (der gesamte Weg) \\
\end{tabular}
\end{merkbox}

\vspace{0.4em}

% ============================================================================
% AUFGABE 1
% ============================================================================
\aufgabe{Aufgabe 1: Setze ein:} \textit{Schrecksekunde, Reaktionsweg, Bremsweg, Anhalteweg}

\vspace{0.3em}
\renewcommand{\arraystretch}{1.6}
\begin{tabularx}{\textwidth}{|X|}
\hline
a) Die Zeit vom Erkennen der Gefahr bis zum Bremsen: \lsg{Schrecksekunde} \\
\hline
b) Der Weg während der Schrecksekunde: \lsg{Reaktionsweg} \\
\hline
c) Der Weg während des Bremsens: \lsg{Bremsweg} \\
\hline
d) Der gesamte Weg bis zum Stillstand: \lsg{Anhalteweg} \\
\hline
\end{tabularx}

\vspace{0.4em}

% ============================================================================
% REAKTIONSZEIT
% ============================================================================
\noindent
\begin{minipage}[t]{0.35\textwidth}
\begin{merkbox}
\centering
\textbf{Meine Reaktionszeit:}\\[0.2em]
{\Large \lsg{ca. 0,5} s}
\end{merkbox}
\end{minipage}
\hfill
\begin{minipage}[t]{0.62\textwidth}
\begin{beispielbox}
\textbf{Was bedeutet ,,10 m/s''?}\\[0.2em]
Du legst \textbf{10 Meter in 1 Sekunde} zurück.\\
In der Schrecksekunde (1\,s) fährst du also 10\,m weiter!
\end{beispielbox}
\end{minipage}

\vspace{0.5em}

% ============================================================================
% AUFGABE 2
% ============================================================================
\aufgabe{Aufgabe 2: Wie weit fährst du in der Schrecksekunde?} \textit{(Überlege: Wie weit in 1 Sekunde?)}

\vspace{0.3em}
\renewcommand{\arraystretch}{1.8}
\begin{tabularx}{\textwidth}{|l|c|X|}
\hline
\textbf{Fahrzeug} & \textbf{Geschwindigkeit} & \textbf{Reaktionsweg} \\
\hline
a) Fahrrad & 5 m/s & In 1 Sekunde: \lsg{5} Meter \\
\hline
b) Auto (Stadt) & 10 m/s & In 1 Sekunde: \lsg{10} Meter \\
\hline
c) Auto (Landstraße) & 20 m/s & In 1 Sekunde: \lsg{20} Meter \\
\hline
\end{tabularx}

\vspace{0.5em}

% ============================================================================
% BREMSWEG-TABELLE
% ============================================================================
\aufgabe{Bremsweg-Tabelle} (für Aufgabe 3 und 4)

\vspace{0.2em}
\begin{merkbox}
\centering
\renewcommand{\arraystretch}{1.3}
\begin{tabular}{|c|c|c|c|}
\hline
\textbf{Geschwindigkeit} & \textbf{Trocken} & \textbf{Nass} & \textbf{Eis} \\
\hline
5 m/s & 2 m & 4 m & 10 m \\
\hline
10 m/s & 5 m & 10 m & 25 m \\
\hline
15 m/s & 10 m & 20 m & 50 m \\
\hline
20 m/s & 20 m & 40 m & 100 m \\
\hline
\end{tabular}
\end{merkbox}

\vspace{0.4em}

% ============================================================================
% AUFGABE 3
% ============================================================================
\aufgabe{Aufgabe 3: Lies den Bremsweg aus der Tabelle ab.}

\vspace{0.2em}
\renewcommand{\arraystretch}{1.6}
\begin{tabular}{|c|c|l|}
\hline
\textbf{Geschwindigkeit} & \textbf{Straße} & \textbf{Bremsweg} \\
\hline
a) 10 m/s & trocken & \lsg{5} m \\
\hline
b) 15 m/s & Eis & \lsg{50} m \\
\hline
c) 20 m/s & nass & \lsg{40} m \\
\hline
\end{tabular}

\newpage

% ============================================================================
% SEITE 2
% ============================================================================

% ============================================================================
% SKIZZE
% ============================================================================
\aufgabe{Skizze: Beschrifte die drei Teile.}

\vspace{0.3em}
\begin{center}
\begin{tikzpicture}[scale=1.0]
    % Straße
    \fill[gray!30] (0,-0.3) rectangle (12,0.3);
    \draw[white, thick, dashed] (0,0) -- (12,0);

    % Auto
    \fill[red!70!black] (0.2,-0.15) rectangle (1.2,0.15);
    \fill[red!50!black] (0.4,0.15) rectangle (1.0,0.4);
    \fill[black] (0.4,-0.22) circle (0.12);
    \fill[black] (1.0,-0.22) circle (0.12);

    % Reaktionsweg (blau) - oberhalb der Straße
    \fill[reaktion!40] (1.4,0.7) rectangle (4.5,1.1);
    \draw[reaktion, thick] (1.4,0.7) rectangle (4.5,1.1);
    \node[loesung, font=\bfseries] at (2.95,1.5) {Reaktionsweg};

    % Bremsweg (orange) - oberhalb der Straße
    \fill[brems!40] (4.5,0.7) rectangle (8.5,1.1);
    \draw[brems, thick] (4.5,0.7) rectangle (8.5,1.1);
    \node[loesung, font=\bfseries] at (6.5,1.5) {Bremsweg};

    % Hindernis
    \fill[red!80!black] (9.5,0) circle (0.35);
    \node at (9.5,0.7) {\small Hindernis};

    % Anhalteweg-Klammer - UNTERHALB der Straße mit Abstand
    \draw[anhalt, line width=1.5pt, decorate, decoration={brace, amplitude=10pt, mirror}]
        (1.4,-0.7) -- (8.5,-0.7);
    \node[loesung, font=\bfseries] at (4.95,-1.5) {Anhalteweg};
\end{tikzpicture}
\end{center}

\vspace{0.4em}

% ============================================================================
% FORMEL
% ============================================================================
\begin{merkbox}
\centering
{\large \textbf{Anhalteweg = Reaktionsweg + Bremsweg}}
\end{merkbox}

\vspace{0.4em}

% ============================================================================
% AUFGABE 4
% ============================================================================
\aufgabe{Aufgabe 4: Berechne den Anhalteweg.} \textit{(Nutze Aufgabe 2 + Tabelle, trockene Straße)}

\vspace{0.3em}
\renewcommand{\arraystretch}{2.0}
\begin{tabularx}{\textwidth}{|l|c|c|c|X|}
\hline
\textbf{Fahrzeug} & \textbf{Geschw.} & \textbf{Reaktionsweg} & \textbf{Bremsweg} & \textbf{Anhalteweg} \\
\hline
a) Fahrrad & 5 m/s & \lsg{5} m & \lsg{2} m & \lsg{7} m \\
\hline
b) Auto (Stadt) & 10 m/s & \lsg{10} m & \lsg{5} m & \lsg{15} m \\
\hline
c) Auto (Landstraße) & 20 m/s & \lsg{20} m & \lsg{20} m & \lsg{40} m \\
\hline
\end{tabularx}

\vspace{0.6em}

% ============================================================================
% BEISPIELRECHNUNG
% ============================================================================
\begin{beispielbox}
\textbf{Beispiel: Auto fährt 10 m/s auf trockener Straße}

\vspace{0.3em}
\begin{tabular}{lll}
\textbf{Reaktionsweg:} & 10 m/s bedeutet 10 m in 1 s & $\rightarrow$ \textbf{10 m} \\
\textbf{Bremsweg:} & aus Tabelle (10 m/s, trocken) & $\rightarrow$ \textbf{5 m} \\
\textbf{Anhalteweg:} & Reaktionsweg + Bremsweg = 10 m + 5 m & $\rightarrow$ \textbf{15 m} \\
\end{tabular}
\end{beispielbox}

\vspace{0.5em}

% ============================================================================
% TRANSFER-AUFGABE (AfB III)
% ============================================================================
\aufgabe{Aufgabe 5: Analysiere die Situation.}

\vspace{0.2em}
\begin{beispielbox}
Ein Ball rollt auf die Straße. Ein Auto fährt mit \textbf{10 m/s} auf \textbf{trockener} Straße.\\
Der Ball ist \textbf{20 Meter} entfernt. Der Fahrer bremst sofort.
\end{beispielbox}

\vspace{0.2em}
a) Berechne den Anhalteweg: \lsg{10} m + \lsg{5} m = \lsg{15} m

\vspace{0.3em}
b) Bewerte: Hält das Auto rechtzeitig vor dem Ball? Begründe mit einer Rechnung.

\vspace{0.4em}
\lsg{Ja, das Auto hält rechtzeitig vor dem Ball.}

\vspace{0.4em}
\lsg{Der Anhalteweg beträgt 15 m, der Ball ist 20 m entfernt. Da 15 m < 20 m, hält das Auto 5 m vor dem Ball.}

% ============================================================================
% EXTENSION 2: Rückwärtsaufgabe
% ============================================================================
\vspace{0.4em}
\aufgabe{Aufgabe 6: Nutze die Tabelle rückwärts!}

\vspace{0.2em}
\begin{beispielbox}
Nach einem Regenschauer sieht man Bremsspuren auf der \textbf{nassen} Straße.\\
Die Bremsspur ist \textbf{20 Meter} lang.
\end{beispielbox}

\vspace{0.2em}
Wie schnell war das Auto?

\vspace{0.4em}
Antwort: Das Auto fuhr mit \lsg{15} m/s

\vspace{0.5em}

% ============================================================================
% SELBSTCHECK
% ============================================================================
\aufgabe{Selbstcheck: Das kann ich jetzt!}

\vspace{0.3em}
\renewcommand{\arraystretch}{1.5}
\begin{tabular}{|p{10cm}|c|c|c|}
\hline
& \textbf{?} & \textbf{OK} & \textbf{!!} \\
\hline
Ich weiß, was die Schrecksekunde ist. & $\square$ & $\square$ & $\square$ \\
\hline
Ich kann den Reaktionsweg bestimmen. & $\square$ & $\square$ & $\square$ \\
\hline
Ich kann den Bremsweg aus der Tabelle ablesen. & $\square$ & $\square$ & $\square$ \\
\hline
Ich kann den Anhalteweg berechnen. & $\square$ & $\square$ & $\square$ \\
\hline
\end{tabular}

\end{document}
